\section{问题三：烘干时长的确定}

\subsection{问题分析与模型的建立}

\subsubsection{本问的判据}

问题三与问题二共用同一套模型，差别只在输出形态。
题目给出的干燥要求是全域各点都达到干燥标准，即\textbf{全域最大干基含水率}
降至 $0.15$~kg/kg 以下：

\begin{equation*}
  \max_{0\le r\le R}C(r,t)\le 0.15~\text{kg/kg}
\end{equation*}

烘干时长即为该条件首次满足的时刻，

\begin{equation*}
  t_{\mathrm{dry}}=\min\left\{t:\ \max_{0\le r\le R}C(r,t)\le 0.15~\text{kg/kg}\right\}
\end{equation*}

\subsubsection{控制点的确定}

需要留意的是判据针对“全域各处”，而由问题二的结果可知，
全域最大干基含水率始终出现在药材中心，故本问以中心点的含水率作为判据的量。

\subsection{结果与分析}

\subsubsection{烘干时长}

\textbf{全域最大干基含水率于 \qthreedryingtime 首次降至判据以下，此即烘干时长。}
表~\ref{tab:table5_q3} 给出题目要求的表格。

\begin{table}[H]
  \centering
  \caption{烘干过程中每隔 $6$~h 的干基水分浓度（kg/kg）}
  \label{tab:table5_q3}
  \small
  \begin{tabular}{lccccc}
    \toprule
    时间/h & $r=0$ & $0.5$~cm & $1.0$~cm & $1.5$~cm & $2.0$~cm（表面） \\
    \midrule
    6 & 1.0170 & 0.9881 & 0.9015 & 0.7550 & 0.5328 \\
    12 & 0.4561 & 0.4432 & 0.4031 & 0.3297 & 0.1642 \\
    18 & 0.2988 & 0.2915 & 0.2687 & 0.2253 & 0.0845 \\
    24 & 0.2378 & 0.2327 & 0.2167 & 0.1855 & 0.0665 \\
    30 & 0.2058 & 0.2018 & 0.1892 & 0.1641 & 0.0598 \\
    36 & 0.1859 & 0.1825 & 0.1718 & 0.1504 & 0.0566 \\
    42 & 0.1721 & 0.1691 & 0.1597 & 0.1407 & 0.0549 \\
    48 & 0.1618 & 0.1592 & 0.1507 & 0.1334 & 0.0537 \\
    54 & 0.1539 & 0.1514 & 0.1436 & 0.1277 & 0.0530 \\
    烘干结束时间 & 0.1500 & 0.1477 & 0.1402 & 0.1249 & 0.0527 \\
    \bottomrule
  \end{tabular}
\end{table}


\subsubsection{穿越判据的单次性}

穿越判据是\textbf{一次性}的，即全域最大干基含水率不会反弹至高于$\le 0.15$~kg/kg。图~\ref{fig:q23} 与图~\ref{fig:q3map} 分别从时程与时空分布两个角度给出这一过程。
前者是中心、$r=1.0$~cm 与表面三个位置的含水率时程，虚线为 $0.15$~kg/kg 判据、点线为烘干时刻。

\begin{figure}[H]
  \centering
  \includegraphics[width=0.98\textwidth]{fig03_drying_curve.pdf}
  \caption{三个位置的含水率时程}
  \label{fig:q23}
\end{figure}

\begin{figure}[H]
  \centering
  \includegraphics[width=0.95\textwidth]{figA2_q3_full_field.pdf}
  \caption{问题三含水率场（附录图）}
  \label{fig:q3map}
\end{figure}

图~\ref{fig:q3map} 把整个烘干过程画成一幅含水率场云图。由该图可以读出两个特征。

\begin{enumerate}[label=\protect\cnum{\arabic*}, leftmargin=2.6em, itemsep=2pt, topsep=3pt, parsep=0pt]
  \item 含水率的下降\textbf{自表面向中心推进}：靠表面的窄带最先变浅，
        中心区在很长时间内仍保持接近初值的高含水率，直到后期才整体变浅。
  \item 越接近烘干结束，整幅云图越趋于均匀，即中心与表面的差别被逐步抹平，
        含水率的下降逐渐减慢，与前文判断的\textbf{干燥由内部水分扩散控制}吻合。
        
\end{enumerate}

\subsection{数值检验与适用性分析}

\subsubsection{数值检验与收敛性}

本问的数值可信度由四类相互独立的证据确认，全部检验量汇总于表~\ref{tab:verify}。

\begin{enumerate}[label=\protect\cnum{\arabic*}, leftmargin=2.6em, itemsep=2pt, topsep=3pt, parsep=0pt]
  \item \textbf{网格加密}：网格由 $400$ 加密到 $800$，报告点上的最大差为 \qthreegridrefinementdiff。
  \item \textbf{时间加密}：内部步长由 $10$~s 减半到 $1$~s，烘干时长的变化为 \qthreetimerefinementdiff。
  \item \textbf{守恒性}：离散质量守恒的相对残差为 \qthreemassbalancerel，达到机器精度量级。
  \item \textbf{独立复算}：格心实现给出的烘干时长与本文方法相差 \qtwodryingtimespread，与问题二一致。
\end{enumerate}

\textbf{四类证据一致表明本问给出的烘干时长已达到网格与步长无关的水平。}
需要指出的是，本问的检验量比问题二更"苛刻"：问题二核对的是 $3$~h 内的温度与含水率，
本问核对的则是判据穿越时刻本身，因此上述差异都落在时间轴上。

\subsubsection{判据尾部平坦性的影响}

烘干时长即为判据的穿越时刻，因此它继承了判据尾部的平坦性：
判据附近中心含水率的下降速率仅约 $8\times10^{-3}$~kg/kg/h，
而本文方法的网格与步长误差在 $10^{-7}$~kg/kg 量级，
落到时间轴上即放大为 \qtwodryingtimespread 的差别，而非方法本身的精度。
故本问给出的烘干时长为 \qthreedryingtime，其时间不确定度约为半分钟，不影响以小时为单位的结论。

\subsubsection{参数与口径灵敏度}

\steptitle{两类被扰动的量}

灵敏度分两类：一是工艺条件量 $h_m$、$D$、$T_\infty$ 与 $C_\infty$；
二是模型自身选定、题目未给的口径，即质量方程的读法。
$h_m$ 取 $\pm5\%$ 以判断响应是否近似线性；
$C_\infty$ 的 $0.045$ 与 $0.055$~kg/kg 是附件 1 稳定段浓度的下限与上限，
两个并列设定都可能代表真实烘房，各算一次；
$D$ 与 $T_\infty$ 只升不降，反向取值不对应任何设问。结果见表~\ref{tab:sens}。

\begin{table}[H]
  \centering
  \caption{单项扰动引起的烘干时长变化}
  \label{tab:sens}
  \small
  \begin{tabular}{lcc}
    \toprule
    扰动项 & 问题二与问题三（相对变化 / \%） & 问题四（绝对变化 / h） \\
    \midrule
    $h_m$ 缩小 $5\%$ & \qtwosenshmminusfivepct & \qfoursenshmminusfivepcthours \\
    $h_m$ 放大 $5\%$ & \qtwosenshmplusfivepct & \qfoursenshmplusfivepcthours \\
    $C_\infty=0.045$~kg/kg & \qtwosenscinflow & \qfoursenscinflowhours \\
    $C_\infty=0.055$~kg/kg & \qtwosenscinfhigh & \qfoursenscinfhighhours \\
    $D$ 放大 $5\%$ & \qtwosensDplusfivepct & \qfoursensDplusfivepcthours \\
    $T_\infty=52$~$^\circ$C & \qtwosensTinfplustwoC & \qfoursensTinfplustwoChours \\
    \bottomrule
  \end{tabular}
  \par\smallskip\footnotesize 正值表示烘干时长变长；$h_m$ 与 $C_\infty$ 各取正负两个方向。
\end{table}

\begin{figure}[H]
  \centering
  \includegraphics[width=0.98\textwidth]{fig12_sensitivity.pdf}
  \caption{烘干时长对单项扰动的响应（问题二与问题三）}
  \label{fig:sens}
\end{figure}

\steptitle{归纳}

表中四项工艺条件的影响一强一弱：$D$ 放大 $5\%$ 与 $T_\infty$ 升高 $2$~$^\circ$C 都缩短时长，属较强一组；
$h_m$ 正负 $5\%$ 两侧量级相同、符号相反，响应近似线性，$C_\infty$ 两个取值一升一降，二者同属最弱一组。
更需声明的是第二类量：主导烘干时长不确定性的是模型口径而非给定参数，
质量方程改用湿基自洽式即变化 \qtwosensmassequationform，量级远超表中任一项；
问题一中表面含水率对 $h_m$ 敏感，其表面值只在题目给定值下成立。

\subsubsection{本问结论的适用范围}

本问的适用范围与问题二相同，另有两点需要说明：

\begin{enumerate}[label=\protect\cnum{\arabic*}, leftmargin=2.6em, itemsep=2pt, topsep=3pt, parsep=0pt]
  \item 恒温段环境取附件 1 稳定段均值外延，属模型设定而非实测；
  \item 一维径向轴对称模型忽略了两端面的轴向输运。
\end{enumerate}